% Generated draft: requires mathematical and source review.
\documentclass[11pt]{article}
\usepackage[margin=25mm]{geometry}
\usepackage{fontspec}
\setmainfont{Latin Modern Roman}
\usepackage{amsmath,amssymb,mathtools,unicode-math}
\setmathfont{Latin Modern Math}
\usepackage{enumitem,needspace,xurl,hyperref}
\setlength{\parindent}{0pt}
\setlength{\parskip}{5pt}
\setlength{\emergencystretch}{3em}
\newcommand{\N}{\mathbb{N}}
\newcommand{\Z}{\mathbb{Z}}
\newcommand{\Q}{\mathbb{Q}}
\newcommand{\R}{\mathbb{R}}
\newcommand{\C}{\mathbb{C}}
\newcommand{\F}{\mathbb{F}}
\newcommand{\A}{\mathbb{A}}
\newcommand{\PP}{\mathbb{P}}
\newcommand{\E}{\mathbb{E}}
\newcommand{\eps}{\varepsilon}
\newcommand{\dd}{\,\mathrm{d}}
\begin{document}
\section*{UnsolvedMath ID: 30006551}
\textbf{Source problem code:} catalogue entry
\subsection{Definitions and mathematical statement}
Let \(q\) be an odd prime power and let \(\F_q\) be the field with \(q\) elements. For a nondegenerate quadratic form \(f\) on a finite-dimensional \(\F_q\)-vector space \(V\), its singular projective points are the one-dimensional subspaces \(\langle v\rangle\) with \(f(v)=0\). A subspace is totally singular if \(f\) vanishes identically on it; the Witt index is the maximum dimension of such a vector subspace. In the standard notation, \(Q^+(2n+1,q)\) arises from a vector space of dimension \(2n+2\) and Witt index \(n+1\), \(Q(2n,q)\) from dimension \(2n+1\) and Witt index \(n\), and \(Q^-(2n+1,q)\) from dimension \(2n+2\) and Witt index \(n\). These definitions apply for positive integers \(n\), but the supplied record does not specify the intended range of \(n\) in the conjecture.

Two distinct singular points are collinear when their span is a two-dimensional totally singular vector subspace. A partial ovoid is a set \(S\) of singular points containing no two distinct collinear points, equivalently meeting each maximal totally singular projective subspace in at most one point. It is maximal if no strictly larger set of singular points is a partial ovoid. Maximality here is with respect to inclusion, not cardinality.

The supplied conjecture says that the listed polar spaces have no maximal partial ovoids of size \(O(q)\). To describe one possible precise asymptotic interpretation, for each fixed \(n\) and each family \(\mathcal P\) among the three families above, put
\[
 m_{\mathcal P,n}(q)=\min\{\lvert S\rvert:S\text{ is a maximal partial ovoid of }\mathcal P_n(q)\}.
\]
If the intended assertion excludes linearly bounded examples along every unbounded sequence of odd prime powers, its fixed-\(n\) formulation would be
\[
 \forall C>0\;\exists q_0\;\forall q\ge q_0\text{ an odd prime power},\qquad m_{\mathcal P,n}(q)>Cq,
\]
or equivalently
\[
 \lim_{\substack{q\to\infty\\q\text{ an odd prime power}}}\frac{m_{\mathcal P,n}(q)}{q}=+\infty.
\]
This is a conditional formulation, not a resolved interpretation of the source. The admissible values of \(n\), whether \(n\) is fixed, the permitted dependence of asymptotic constants on \(n\) and the family, and whether bounded examples on subsequences are excluded are not specified. Merely excluding a single \(O(q)\)-bounded construction valid for every sufficiently large odd prime power would be a different assertion. These qualifications require clarification before an unconditional quantified statement can be assigned to the source.
\subsection{Short English statement}
The supplied conjecture asks whether maximal partial ovoids in the three listed families of quadrics must grow faster than linearly in the odd field order \(q\). A maximal partial ovoid is an inclusion-maximal set of mutually noncollinear singular points. The intended dimension range and asymptotic quantifiers remain unresolved.
\subsection{Sources}
\begin{itemize}
\item Local UnsolvedMath registry: supplied complete record for permanent ID \(30006551\), title \emph{Small Maximal Partial Ovoids in Polar Spaces}, exact source problem code \texttt{catalogue entry}; Oberwolfach Reports Open Problems collection, proposed year \(2026\). No individual proposer is identified.
\item Supplied report citation: \emph{Combinatorics} (2026), Oberwolfach Reports, DOI 10.4171/owr/2026/1, pp. 5--79. \url{https://doi.org/10.4171/owr/2026/1}. The record describes these reports as documenting mathematical workshops and their open-problem discussions.
\item Supplied literature reference: Schillewaert--Verstraete, \emph{Small maximal partial ovoids}. \url{https://arxiv.org/abs/1308.1899}.
\item The original excerpt attributes the conjecture to reference [13], whose bibliographic entry is not supplied. Its acknowledgments state that the work was joint in parts with Chase Wilson, Jiaxi Nie and Jeroen Schillewaert. These acknowledgments do not identify an individual proposer of this conjecture.
\item Supplied literature assessment, checked on exactly \texttt{2026-08-22}: status \emph{open}; the record states that the \(O(q)\) lower-scale conjecture is not proved and that existing constructions give \(O(q\log q)\) or comparable polylogarithmic bounds. Both linked references above are listed as sources checked. This is supplied dated curation derived from the Oberwolfach record, not an independently verified literature assessment; its status and cited claims require independent checking.
\item The registry labels the statement \texttt{corrected\_verified} and reports that the official report was checked, the question isolated from surrounding exposition, and notation repaired for self-contained publication. These are supplied provenance claims, not verification performed for this draft.
\end{itemize}
\end{document}
